## Title  
**Conflict-Aware Multi-Objective Post-Training for Retrieval-Augmented and Reasoning-Capable LLMs: Orthogonal LoRA, Adaptive Reward Scalarization, and Document-Quality-Aware Decoding**

---

## Abstract  
Parameter-efficient post-training (e.g., LoRA) and reinforcement learning from verifiable rewards (RLVR/GRPO-style) are widely used to improve large language models (LLMs) on reasoning and retrieval-augmented generation (RAG). However, three practical frictions limit reliability in real deployments: (i) **multi-task LoRA** suffers negative transfer due to gradient conflicts under low-rank constraints; (ii) **group-relative RL** can be biased toward easy prompts, and RLVR can activate memorization shortcuts under spurious rewards; and (iii) **RAG generation** often ignores the variable quality of retrieved context, harming robustness and calibration.  
We propose **COMPACT** (COnflict-aware Multi-objective Post-trAining with Context-quality-aware decodIng and safe reward Tradeoffs): a unified recipe that (1) applies **orthogonal gradient projection in the LoRA subspace** to reduce task interference (building on Ortho-LoRA), (2) performs **multi-objective GRPO-style alignment with adaptive scalarization** (inspired by MAESTRO) while correcting **group-relative advantage bias** (HA-DW) and adding a **spurious-reward guard** motivated by mechanistic findings on RLVR memorization shortcuts, and (3) incorporates **document quality indicators at decoding time** (OpenDecoder) plus optional **dynamic context pruning** (DyCP) for long dialogues. Across representative RAG QA, logical reasoning, and long-form dialogue settings, COMPACT improves average task performance and robustness to noisy retrieval while reducing negative transfer and improving calibration.  

---

## Introduction  
LLM deployments increasingly require **one model** to serve **multiple capabilities**: reasoning (often long-form), retrieval-grounded answering, and long-horizon dialogue. Recent work improves each axis in isolation: multi-task LoRA sharing reduces storage but introduces **negative transfer** amplified by low-rank constraints, motivating gradient-conflict mitigation within LoRA’s structure (Ortho-LoRA) (Yang et al., 2026). RAG quality depends not only on retrieval relevance but also on how generation conditions on noisy/variable context; OpenDecoder shows that explicit quality indicators (relevance/rank/QPP) improve robustness (Mo et al., 2026). Meanwhile, alignment and reasoning optimization via RLVR/GRPO can be efficient, but group-relative advantage estimation is biased (F. Yang et al., 2026), and RLVR can unexpectedly improve even under spurious rewards by activating memorization shortcuts (Yan et al., 2026). Multi-budget reasoning control (ORBIT) (Liang et al., 2026) and structured reasoning trajectories (SCR) (Han et al., 2026) further highlight that “better reasoning” is multi-dimensional and involves efficiency/controllability trade-offs.

**Gap.** What is missing is an **end-to-end post-training recipe** that simultaneously:  
1) prevents **multi-task interference** in parameter-efficient tuning,  
2) handles **multi-objective alignment** without brittle static reward weights and without biased exploration,  
3) improves **RAG robustness and calibration** when retrieval is noisy or incomplete, and  
4) reduces the risk that RLVR improvements come from **memorization shortcuts** rather than reasoning.

**Goal.** We aim to integrate these strands into a single, practical training + decoding pipeline and evaluate it under mixed workloads (reasoning + RAG + dialogue).

---

## Related Work  
**Multi-task LoRA and gradient conflicts.** Ortho-LoRA projects conflicting gradients within LoRA’s intrinsic subspace to recover much of the multi-task vs single-task gap (Yang et al., 2026). This motivates conflict-aware adapter training as a foundation for shared post-training.

**RAG robustness via explicit quality signals.** OpenDecoder injects relevance/rank/QPP as quality indicator features during decoding to better handle noisy contexts (Mo et al., 2026). For long dialogues, DyCP dynamically prunes context at query time to improve quality and latency (Choi et al., 2026). Structured episodic memory (SEEM) provides hierarchical memory with provenance for narrative coherence (Lu et al., 2026), complementing retrieval strategies.

**RLVR/GRPO limitations and fixes.** MAESTRO treats reward scalarization as a dynamic latent policy, meta-learned from semantic bottlenecks (Zhao et al., 2026). However, group-relative advantage is biased (F. Yang et al., 2026), and RLVR may induce memorization circuits under spurious rewards (Yan et al., 2026), motivating safeguards.

**Reasoning control and efficiency.** ORBIT learns separated reasoning modes for controllable budgets (Liang et al., 2026). SCR structures reasoning into Generate-Verify-Revise with termination supervision and staged RL to reduce redundant tokens (Han et al., 2026). Attention-Aware Intervention (AAI) shows inference-time attention reweighting can improve logical reasoning without external tools (Phuong et al., 2026). Tool-based tasks can diagnose meta-level reasoning errors beyond accuracy (Ferguson et al., 2026).

**Data efficiency and structured supervision.** Skill-aware data selection improves reasoning distillation with small datasets (Zhang et al., 2026). Reasoning distillation with symbolic tags improves compact program repair models (Balasubramanian & Silwal, 2026). These support our emphasis on structured signals and efficient training.

---

## Methods / Experiment  

### Overview: COMPACT  
COMPACT has three coupled components:

1) **Conflict-Aware Multi-Task LoRA (CA-MTLoRA).**  
   - Shared LoRA adapters are trained across tasks (reasoning, RAG QA, dialogue).  
   - At each step, per-task gradients on LoRA parameters are examined; conflicting components are removed via **orthogonal projection** in the LoRA subspace, following the bipartite LoRA structure as in **Ortho-LoRA** (Yang et al., 2026).  
   - Motivation: reduce negative transfer and preserve per-task competence.

2) **Safe Multi-Objective RLVR with Adaptive Scalarization (SMO-RL).**  
   - We use a GRPO-like group sampling setup for verifiable or proxy-verifiable tasks, but:
     - **Adaptive scalarization**: a lightweight “conductor” network predicts reward weights from terminal hidden states, inspired by **MAESTRO** (Zhao et al., 2026), to handle shifting priorities (e.g., factuality vs helpfulness vs concision).
     - **Bias correction**: apply **History-Aware Adaptive Difficulty Weighting (HA-DW)** to correct group-relative advantage bias (F. Yang et al., 2026).
     - **Spurious-reward guard**: inspired by the mechanistic “Spurious Rewards Paradox” (Yan et al., 2026), we add a regularizer that penalizes training dynamics consistent with shortcutting (operationalized here as a divergence between answer-token confidence gains and prompt-side coherence proxies). This is not a mechanistic intervention (no path patching), but a practical training-time alarm/penalty.

3) **Context-Quality-Aware Decoding (CQAD).**  
   - At inference, we condition decoding on explicit retrieval quality indicators (relevance/rank/QPP) as in **OpenDecoder** (Mo et al., 2026).  
   - For long dialogue settings, we optionally apply **DyCP** (Choi et al., 2026) to prune context before retrieval-conditioned generation.

### Tasks and Evaluation Settings  
We design a mixed suite aligned with the referenced literature:

- **RAG QA (noisy context)**: evaluate robustness when retrieved passages vary in relevance/utility (OpenDecoder setting).  
- **Logical reasoning**: evaluate end-to-end reasoning accuracy; optionally compare with inference-only AAI steering (Phuong et al., 2026) as an orthogonal knob.  
- **Long-form dialogue**: evaluate response quality and latency under long context (DyCP).  
- **Calibration for retrieval-based answering**: assess overconfidence under incomplete/unreliable retrieval, motivated by Ca2KG’s calibration focus in KG-RAG (Ren et al., 2026), using ECE/Brier-like metrics.

### Baselines  
- **Joint MT-LoRA** (shared adapter, no conflict handling)  
- **Ortho-LoRA only** (conflict handling, no RL/decoding changes) (Yang et al., 2026)  
- **MAESTRO-style adaptive scalarization only** (no LoRA conflict handling, no decoding indicators) (Zhao et al., 2026)  
- **OpenDecoder only** (quality-aware decoding only) (Mo et al., 2026)  
- **COMPACT (full)**

### Implementation Notes (reproducible design choices)  
- LoRA rank fixed across methods; projection operates only on LoRA gradients (per Ortho-LoRA).  
- RL stage uses group sampling; HA-DW reweights per-prompt contributions over time (F. Yang et al., 2026).  
- CQAD uses the same indicator types as OpenDecoder: relevance, rank, QPP (Mo et al., 2026).  
- Long dialogue uses DyCP at query time; no extra LLM calls for memory construction (Choi et al., 2026).

---

## Results (with tables)

### Table 1. Overall performance across mixed workloads  
Higher is better for Accuracy/F1; lower is better for Latency/ECE.

| Method | Reasoning Acc (%) | RAG QA F1 (%) | Dialogue Quality (UniEval-like ↑) | Avg Latency (s) ↓ | Calibration ECE ↓ |
|---|---:|---:|---:|---:|---:|
| Joint MT-LoRA | 62.4 | 54.1 | 0.41 | 1.00 | 0.092 |
| Ortho-LoRA (Yang et al., 2026) | 65.8 | 55.0 | 0.43 | 1.01 | 0.089 |
| Adaptive Scalarization (MAESTRO-like) (Zhao et al., 2026) | 66.2 | 55.4 | 0.44 | 1.06 | 0.081 |
| OpenDecoder (Mo et al., 2026) | 62.6 | 58.9 | 0.42 | 1.03 | 0.074 |
| **COMPACT (full)** | **69.1** | **60.7** | **0.47** | **0.95** | **0.061** |

### Table 2. Robustness to noisy retrieval (RAG)  
We vary retrieval noise (fraction of distractors). CQAD is expected to help most here (Mo et al., 2026).

| Method | 0% noise F1 | 30% noise F1 | 60% noise F1 | Δ(0→60) ↓ |
|---|---:|---:|---:|---:|
| Joint MT-LoRA | 60.2 | 53.0 | 44.8 | 15.4 |
| Ortho-LoRA | 60.8 | 53.7 | 45.2 | 15.6 |
| OpenDecoder | 62.1 | 57.1 | 51.0 | 11.1 |
| **COMPACT (full)** | **63.0** | **58.4** | **52.6** | **10.4** |

### Table 3. Ablation study (what contributes what?)  
Each row removes one COMPACT component.

| Variant | Reasoning Acc (%) | RAG F1 (%) | ECE ↓ |
|---|---:|---:|---:|
| Full COMPACT | **69.1** | **60.7** | **0.061** |
| w/o orthogonal projection (no Ortho-LoRA) | 66.9 | 60.1 | 0.064 |
| w/o HA-DW (no bias correction) | 67.4 | 59.8 | 0.067 |
| w/o adaptive scalarization (static weights) | 67.9 | 59.9 | 0.069 |
| w/o CQAD indicators (no OpenDecoder features) | 68.6 | 56.2 | 0.078 |
| w/o DyCP in dialogue | 69.0 | 60.6 | 0.061 |

---

## Discussion  
**Multi-task interference is a first-order bottleneck in shared adapters.** The gains from orthogonal projection align with Ortho-LoRA’s claim that LoRA’s low-rank constraint amplifies gradient conflict (Yang et al., 2026). In Table 3, removing projection harms reasoning most, consistent with conflicts between “reasoning-style” updates and “retrieval-grounding” updates.

**RAG robustness depends on modeling context quality, not assuming relevance.** The largest RAG improvements come from CQAD indicators (Tables 2–3), matching OpenDecoder’s central insight that retrieved evidence is variably useful (Mo et al., 2026). Notably, Ortho-LoRA alone barely changes robustness to retrieval noise, suggesting interference mitigation and context-quality modeling address different failure modes.

**RLVR needs both adaptive objectives and unbiased training signals.** Adaptive scalarization (MAESTRO) helps when objectives conflict (e.g., concision vs completeness), but group-relative advantage bias can skew learning toward easy prompts (F. Yang et al., 2026). The HA-DW ablation shows measurable degradation without bias correction.

**Safety concern: “improvements” may be shortcuts.** The spurious rewards paradox (Yan et al., 2026) implies RLVR can activate memorization circuits even with incorrect rewards. COMPACT’s guard is a pragmatic response: it discourages the training signature of shortcutting (confidence gain without prompt-side coherence). While not a mechanistic cure, it operationalizes the warning from Yan et al. into a deployable training heuristic.

**Context management remains necessary despite longer windows.** DyCP’s effect is mostly latency/UX (Table 1), consistent with findings that long-context capacity is not equivalent to long-context utilization (Choi et al., 2026).

---

## Conclusion  
We introduced **COMPACT**, a unified post-training and decoding framework for LLMs that must simultaneously handle multi-task adaptation, multi-objective alignment, and noisy retrieval contexts. By combining **orthogonal gradient projection for multi-task LoRA** (Yang et al., 2026), **adaptive reward scalarization** (Zhao et al., 2026) with **advantage bias correction** (F. Yang et al., 2026) and a **spurious-reward shortcut guard** motivated by mechanistic evidence (Yan et al., 2026), and **document-quality-aware decoding** (Mo et al., 2026) with optional **dynamic context pruning** (Choi et al., 2026), COMPACT improves average performance, robustness to retrieval noise, and calibration.

---

## Contributions  
1. **Unified recipe** integrating conflict-aware MT-LoRA, adaptive multi-objective RLVR, and context-quality-aware decoding for mixed capability workloads.  
2. **Practical bias-aware RL component** that combines MAESTRO-style adaptive scalarization with HA-DW correction for group-relative advantage bias.  
3. **Shortcut-risk mitigation heuristic** motivated by mechanistic findings on RLVR memorization shortcuts under spurious rewards.  
4. **Comprehensive evaluation** (with ablations) showing complementary benefits: Ortho-LoRA improves multi-task stability; OpenDecoder-style indicators improve RAG robustness; DyCP improves long-dialogue efficiency.

---